\section{Disposición de los diagramas}
\label{subsec:disposicion}

Varios tableros pueden ir en una grilla: los |goban| y los diagramas de |\sgfdiagrams| y |\sgfproblem|, dentro del entorno |pgfgolayout|, o los de un solo |\sgfdiagrams| o |\sgfproblem|, con su \emph{key} |layout|. Sin grilla, los diagramas de esos comandos van en el párrafo, uno al lado del otro, y pasan a la línea siguiente cuando no entran. Con |columns|, van en una grilla de columnas del mismo ancho: cada fila es una línea aparte (la página puede cortarse entre filas), los tableros de una fila quedan alineados arriba y una fila incompleta deja vacías las últimas columnas. Con |fit|, además, los tableros se escalan para llenar el ancho de su columna.

\begin{environment}{{pgfgolayout}\opt{\oarg{layout options}}}
	Dispone en una grilla los |goban| que contiene, y los diagramas de los |\sgfdiagrams| y |\sgfproblem| que contiene, en el orden en que aparecen. Las opciones son las de |/pgfgo/layout|; las opciones |layout| de los comandos de adentro no se usan. Adentro no debe haber texto, solo los tableros y los comandos. Con |fit|, cada |goban| se dibuja varias veces para medirlo: su contenido no debe tener efectos fuera del tablero (por ejemplo, contadores que aumentan).

	\begin{codeexample}[code only]
		\begin{pgfgolayout}[columns = 4, fit]
			\begin{goban}[5]\stones[b]{C3}\end{goban}
			\begin{goban}[5]\stones[w]{C3}\end{goban}
			\sgfdiagrams[step = 5]{partida}
		\end{pgfgolayout}
	\end{codeexample}
\end{environment}

\begin{pgfgolayout}[columns = 4, fit]
	\begin{goban}[5]\stones[b]{C3}\end{goban}
	\begin{goban}[5]\stones[w]{C3}\end{goban}
	\sgfdiagrams[step = 5]{partida}
\end{pgfgolayout}

\begin{codeexample}[code only]
	\sgfdiagrams[step = 2, to = 6, layout = {columns = 3, fit}]{partida}
\end{codeexample}

\sgfdiagrams[step = 2, to = 6, layout = {columns = 3, fit}]{partida}

\begin{stylekey}{/pgfgo/layout = \meta{layout options}}
	Opciones de la disposición para todo el documento (por ejemplo, |columns = 3, fit|). Para un solo comando, |/pgfgo/sgf diagrams/layout| y |/pgfgo/sgf problem/layout|; también para todo el documento, |\pgfgoset{sgf diagrams = {layout = {...}}}|.
\end{stylekey}

\begin{key}{/pgfgo/layout/columns = \meta{number o auto} (initially auto)}
	Cantidad de columnas de la grilla. Con |auto|, los diagramas van en el párrafo, sin grilla.
\end{key}

\begin{key}{/pgfgo/layout/width = \meta{dimension} o |natural| (initially \texttt{\string\linewidth})}
	Ancho de la grilla. Si es menor que el de la línea, la grilla se centra. Cada columna mide lo que queda al restarle los |column sep| a |width|, repartido entre las columnas; un tablero más angosto queda en su columna según |halign|, así que entre dos tableros hay más espacio que |column sep|.

	Con |natural|, la grilla se ajusta a los diagramas: las columnas tienen el ancho del diagrama más ancho (con su epígrafe, que se corta como sin grilla, ver |caption/min width|) y entre dos tableros del mismo ancho queda exactamente |column sep|. La grilla mide eso por la cantidad de columnas (o de diagramas, si son menos) y se centra en la línea; si es más ancha que la línea, se sale del margen. Los diagramas se dibujan todos juntos al final (o al terminar cada grupo de |per page| o |per frame|). |fit| no sirve con |natural|: el paquete avisa y los tableros quedan con su escala.

	\begin{codeexample}[code only]
		\begin{pgfgolayout}[columns = 2, width = natural, column sep = 0em, row sep = 0em]
			\begin{goban}[partial = A9:E5, scale = 0.5]\stone[b]{C7}\end{goban}
			\begin{goban}[partial = E9:I5, scale = 0.5]\stone[w]{G7}\end{goban}
			\begin{goban}[partial = A1:E5, scale = 0.5]\stone[b]{C3}\end{goban}
			\begin{goban}[partial = E1:I5, scale = 0.5]\stone[w]{G3}\end{goban}
		\end{pgfgolayout}
	\end{codeexample}

	\begin{pgfgolayout}[columns = 2, width = natural, column sep = 0em, row sep = 0em]
		\begin{goban}[partial = A9:E5, scale = 0.5]\stone[b]{C7}\end{goban}
		\begin{goban}[partial = E9:I5, scale = 0.5]\stone[w]{G7}\end{goban}
		\begin{goban}[partial = A1:E5, scale = 0.5]\stone[b]{C3}\end{goban}
		\begin{goban}[partial = E1:I5, scale = 0.5]\stone[w]{G3}\end{goban}
	\end{pgfgolayout}
\end{key}

\begin{key}{/pgfgo/layout/column sep = \meta{dimension} (initially 1em)}
	Espacio entre columnas.
\end{key}

\begin{key}{/pgfgo/layout/row sep = \meta{skip} (initially \texttt{\string\medskipamount})}
	Espacio entre filas.
\end{key}

\begin{key}{/pgfgo/layout/fit = \meta{boolean} (default true, initially false)}
	Escala cada tablero para que llene el ancho de su columna; reemplaza el |scale| del tablero. Las líneas no cambian de grosor. En |\sgfdiagrams| y |\sgfproblem|, todos los tableros usan la escala que llena la columna con el primero, así las piedras tienen el mismo tamaño en todos. Si con esa escala un tablero no entra en la columna (un problema más ancho que el primero, por ejemplo), solo ese tablero se achica hasta llenarla. Solo con |columns|, y no con |width = natural|.
\end{key}

\begin{key}{/pgfgo/layout/per page = \meta{number} o |false| (initially 0)}
	Cada \meta{number} diagramas, empieza una página nueva (con |\newpage|): con |per page = 1|, un diagrama por página. En la grilla, una fila a medias se dibuja antes del corte; conviene que \meta{number} sea un múltiplo de |columns|. Con |0| o |false|, no hay cortes. En |\sgfproblem| con |statement = alone|, el enunciado cuenta como un diagrama.

	\begin{codeexample}[code only]
		\sgfdiagrams[step = 50, layout = {columns = 2, fit, per page = 4}]{lee}
	\end{codeexample}
\end{key}

\begin{key}{/pgfgo/layout/per frame = \meta{number} o |false| (default 1, initially 0)}
	Dentro de un |frame| de |beamer|, cada \meta{number} diagramas pasan a la diapositiva siguiente: con |per frame| sola, un diagrama por diapositiva, cada uno con su epígrafe. Los diagramas usan |\only<+>|, así que se suman a las demás capas del |frame|: después de un |\pause|, el primer diagrama va en la diapositiva siguiente. Fuera de |beamer|, el paquete avisa y la \emph{key} no se usa. En |beamer|, si también hay |per page|, vale |per frame|.

	\begin{codeexample}[code only]
		\begin{frame}{La partida}
			\sgfdiagrams[step = 50, goban = {scale = 0.6}, layout = {per frame}]{lee}
		\end{frame}
	\end{codeexample}
\end{key}

\begin{key}{/pgfgo/layout/halign = \meta{center, left o right} (initially center)}
	Dónde va cada diagrama dentro de su columna, con su epígrafe y sus notas. Los comentarios van siempre alineados a la izquierda.

	\begin{codeexample}[code only]
		\sgfdiagrams[step = 2, to = 6, layout = {columns = 2, width = 0.6\linewidth,
			column sep = 3em, row sep = 2em, halign = left}]{partida}
	\end{codeexample}

	\sgfdiagrams[step = 2, to = 6, layout = {columns = 2, width = 0.6\linewidth,
		column sep = 3em, row sep = 2em, halign = left}]{partida}
\end{key}

\begin{key}{/pgfgo/layout/valign = \meta{top o bottom} (initially top)}
	Cómo se alinean los diagramas de una fila de la grilla cuando sus tableros tienen distinto alto. Con |top|, arriba. Con |bottom|, los tableros terminan a la misma altura, y sus epígrafes y notas empiezan todos en la misma línea, debajo. Con |columns = auto|, los diagramas van en la línea de base del párrafo, como con |bottom|.

	\begin{codeexample}[]
		\begin{pgfgolayout}[columns = 3, valign = bottom]
			\begin{goban}[5, scale = 0.6]\stones[b]{C3}\end{goban}
			\begin{goban}[5x3, scale = 0.6]\stones[w]{C2}\end{goban}
			\begin{goban}[9, scale = 0.6]\stones[b]{E5}\end{goban}
		\end{pgfgolayout}
	\end{codeexample}
\end{key}

\begin{key}{/pgfgo/layout/align captions = \meta{boolean} (default true, initially false)}
	Con |valign = top|, los epígrafes, las notas y los comentarios de los diagramas de una fila empiezan en la misma línea, debajo del tablero más alto; los tableros siguen alineados arriba. Con |valign = bottom| ya están en la misma línea. Los |goban| de |pgfgolayout| no tienen epígrafe: no cambian.

	\begin{codeexample}[code only]
		\sgfproblem[game = 1:3, solutions = false, partial = auto,
			layout = {columns = 3, align captions}]{problemas}
	\end{codeexample}
\end{key}
